\documentclass{ximera}
\title{Some Basic Trigonometric Identities}

\newcommand{\pskip}{\vskip 0.1 in}

\begin{document}
\begin{abstract}
Using the sine and cosine functions to describe simple harmonic motion.
\end{abstract}
\maketitle


\begin{question} \label{Qppf3}
Use the Pythagorean theorem to find the length of the altitude of an equilateral triangle with side length $2$ cm. What about if the equilateral triangle has side length $74$cm?
\end{question}

\begin{question} \label{Qof3efr3}
\begin{enumerate}
\item An isoceles right triangle has legs of length $1$ cm. Find the length of the hypotenuse.

\item An isoceles right triangle has legs of length $14$ cm. Find the length of the hypotenuse.

\item An isoceles right triangle has a hypotenuse with length $10$ cm. Find the length of the legs.
\end{enumerate}
\end{question}


\begin{question} \label{Qpefere3}
Work on a circle of radius $8$cm to find the exact values of each of the following. Draw an equilateral triangle or an isosceles right triangle in your coordinate system as appropriate.

\begin{enumerate}
\item $\cos(7\pi/6)$

\item $\sin(83\pi/2)$

\item $\cos(11\pi)$

\item $\sin(11\pi/6)$

\item $\cos(-3\pi/4)$
\end{enumerate}
\end{question}

\begin{question} \label{Qpfrfsfd}
Let $P$ be the point with coordinates $(4,1)$ and let $\theta$ be a polar angle of the point $P$. Use geometry to find the exact values of each of the following. Draw pictures to help with your explanations. Label all angles and coordinates. Use only geometry and the definitions of $\sin\theta$ and $\cos\theta$. No formulas.

\begin{enumerate}
\item $\sin(\theta + \pi/2)$

\item $\sin(\theta + 3\pi/2)$

\item $\cos (\theta + \pi)$

\item $\sin(\theta - 6\pi)$

\item $\cos(\theta + \pi/2)$

\item $\cos(-\theta)$

\item $\sin(-\theta)$

\item $\cos(\theta + \pi/4)$
\end{enumerate}
\end{question}

\begin{question} \label{Qdfefr30}
Use the ideas of the previous question to express each of the following in terms of $\cos\theta$ and/or $\sin\theta$. Draw a picture to help with your explanation. Start by drawing the point $P$ with coordinates $(x,y)$ and the circle with radius $r=\sqrt{x^2 + y^2}$ centered at the origin through $P$. Label all angles and coordinates. Use only geometry and the definitions of $\sin\theta$ and $\cos\theta$. No formulas. 

\begin{enumerate}
\item $\sin(\theta + \pi/2)$

\item $\sin(\theta + 3\pi/2)$

\item $\cos (\theta + \pi)$

\item $\sin(\theta - 6\pi)$

\item $\cos(\theta + \pi/2)$

\item $\cos(-\theta)$

\item $\sin(-\theta)$

\item $\cos(\theta + \pi/4)$
\end{enumerate}
\end{question}

\end{document}
